\scanpage{84}
Soit $I_{\widetilde x}$ le noyau de l'hom. obtenu (ou stabilisateur ``géométrique'' de $\widetilde x$), donc on a une suite exacte
\[
(7)\quad 1-I_{\widetilde x}\longrightarrow N_{\widetilde x}\longrightarrow\operatorname{Gal}(\overline{k(x)}/k(x))\longrightarrow1
\]
et par Kummer un iso. can.
\[
(8)\quad I_{\widetilde x}\simeq T_\infty(\overline{k(x)}^*)\simeq T_\infty(\overline K^*)
\]
\margin{à cause de l'hyp de car. 0 !}
On notera que si $x$ est le pt gén. du diviseur $D$, alors $k(x)$ est le corps des fonctions de $D$. C'est un corps de type fini sur $\mathbb Q$.

Il est immédiat (sans supposer que le corps de base pour $S$ soit $\mathbb Q$) que le noyau de $\deleted{\unsure{\pi|\Gamma}}\inserted{E}_K\longrightarrow E_S$ est le sous-groupe invariant engendré par les $I_{\widetilde x}$. Donc l'hypothèse que $E_K\longrightarrow E_{U_S}$ s'annule sur ledit noyau, signifie aussi qu'il s'annule sur chacun des $I_{\widetilde x}$.

Soit \unsure{donc|sur} $U_{\underline O_x}$ induit par $U$ sur $\operatorname{Spec}O_x$, on a donc des factorisations
\margin{83}
